\section{Conclusion}

Expanding the \YazSes{} evaluation beyond one clean-speech WER changed several engineering
conclusions.

Harder speech altered the relative robustness of local recognizers. Repeated decoding exposed a
\texttt{large-v3} failure mode in which insertion count and runaway continuation vary while
substitution, deletion, and correct-word counts remain stable. Previous-text conditioning
behaved differently across checkpoints rather than offering one global rule. Cross-platform
decoding showed model-dependent numerical variation. Synthetic onset padding did not recover
missed speech, and CPU streaming was beneficial only under a sufficiently fast checkpoint and
compute budget. In Meeting Mode, real annotated AMI data exposed a catastrophic mismatch in the
former clustering operating point, while known speaker count proved useful for count correctness
without establishing a DER improvement.

The broader lesson is methodological. Offline interactive speech systems should not be evaluated
only by a single average accuracy number. Reproducibility, tail failures, hardware dependence,
latency under the actual interactive workload, and task-specific metrics can change both the
scientific interpretation and the product default.

Equally important, contradictory measurements should be preserved rather than overwritten. In
this campaign, retaining disagreement was not bookkeeping; it created the experiment that
revealed one of the paper's central findings.
